	\subsection{同类猜想} \label{ss:variants}
	本节讨论了局部几何朗兰兹猜想，希望读者对其陈述、以及它在表示论与
	几何中的一些具体预测（既有已知的，也有猜想的）有了基本的感受。

	在结束之前，我们应当提到附近领域中一些重要的变体，或者更确切地说，同类的猜想，以及
	一些其他正在浮现的联系。

	\ssszh[拟分裂群] 首先要作的评注是：本综述从头到尾使用的都是环路群，即 $k(\!(t)\!)$ 上的分裂
	约化群。然而，即便为了研究它们的范畴表示，正如数论中那样，考虑扭环路群的情形也很重要，
	即 $p$-adic 情形中拟分裂群的范畴类比。注意这也为扭仿射李代数的表示论提供了类似的图景，而
	下面讨论的局部几何朗兰兹对偶的各种形式，预期同样可以推广到拟分裂情形。

	\ssszh[量子局部几何朗兰兹] 回忆 $\fgk\mod$
	只是对整水平才是环路群的范畴表示。对其他水平，它反而带有扭 D-模范畴 $\Dmod_\kappa(LG)$ 的作用。
	可以问：对环路群的这类扭范畴表示，是否存在某种形式的局部朗兰兹对偶。确实存在，这正是量子几何
	朗兰兹的主题。

    为继续下去，为简单起见假设
    $\kappa - \kappa_c$ 是非退化双线性型。此时可以给 $\cG$ 配上一个对偶型 $\check{\kappa}$，如下。若以
    $\kappa_{\cfg, c}$ 记 $\cfg$ 的临界水平，则 $\check{\kappa}$ 由下述性质定义：
    \[
\kappa - \kappa_c \text{ 与 } \check{\kappa} - \kappa_{\cfg, c}
\]
    分别是 $\ft$ 与 $\check{\ft}$ 上的对偶双线性型。也就是说，在差一个我们反复遇到的临界平移之后，
    要求 $\kappa$ 与 $\check{\kappa}$ 互为对偶型。

    由 Gaitsgory \cite{quantum-langlands-summary} 提出的局部量子朗兰兹猜想是：存在 $(\infty,2)$-范畴的
    等价
    \[
\Dmod_\kappa(LG)\mod \simeq \Dmod_{-\check{\kappa}}(L\cG)\mod.
\]
    这建立在整体情形中 Beilinson--Drinfeld、Feigin--Frenkel 与 Stoyanovsky \cite{stoyanovsky} 的
    早期工作之上。

     这里只作几点评注；更多细节可参见 \cite{winter-langlands-summary}、\cite{paris-notes}。

    \begin{re} 关于量子情形，一个基本要点是：两端看起来更对称，因为它们只涉及扭的范畴环路群
    表示。此外，在确切的意义下，当 $\kappa \rightarrow \kappa_c$、从而 $\check{\kappa} \rightarrow \infty$ 时，
    就恢复通常的陈述 \eqref{e:loclang}；而当 $\check{\kappa} \rightarrow \kappa_{\cfg, c}$、从而
    $\kappa \rightarrow \infty$ 时，就恢复把 $G$ 与 $\cG$ 的角色互换后的类比。
    	\end{re}

    \begin{re} 「量子」这一形容词的来源如下。回忆 Feigin--Frenkel 同构
    %
    \[      \mathscr{W}_{\fg, \kappa_c} \simeq \OO_{\Op},
    \]
    %
    它在临界水平的故事中起基本作用。事实上，Feigin--Frenkel 在 \cite{feiginfrenkel91} 中证明了：
    它是仿射 W-代数的一族同构
    %
    \[   \mathscr{W}_{\fg, \kappa} \simeq \mathscr{W}_{\cfg, \check{\kappa}} \]
    %
    的准经典极限。这些同构（如今是非交换代数之间的同构）是前一个同构的形变量子化，在量子理论中
    起基本作用。

    与此相关，在临界水平我们看到对偶群 $\on{Rep}(\cG)$ 的表示在几何朗兰兹的非分歧情形中起基本作用；
    而在量子设定下，起类似作用的是量子群 $\on{Rep}_q(\cG)$ 的表示。
    \end{re}

   \ssszh[与物理的联系] 迄今为止，我们只触及了几何朗兰兹与物理之间的一些关系，而且大多是与二维
   共形场论的联系。然而，它与三、四维超对称规范理论之间有极其丰富的联系，这始于
   \cite{kapustinwitten}、\cite{wittenwild}、\cite{gukovwitten}、\cite{frenkelwitten} 等文。这是（量子）几何朗兰兹中
   许多有趣构造与预期相容性的来源，无论局部还是整体。部分近期工作可参见
   \cite{gaiotto16s-duality}、\cite{gaiottorapcak}、\cite{bfn}、\cite{gaiottotwistcompact}、\cite{costellogaiotto}、
   \cite{creutziggaiotto}、\cite{creutziggaiottolinshaw}、\cite{frenkelgaiotto}、\cite{bfgt}、\cite{hilburnraskin}，
   以及即将发表的 \cite{gaiottoyu}。

 \begin{re}
   上述部分工作的一个令人兴奋之处，是它与算术朗兰兹纲领的新进展密切相关。也就是说，由四维
   $\mathscr{N} = 4$ Yang--Mills 的边界条件的 S-对偶给出的朗兰兹对偶范畴的许多预测
   （Gaiotto--Witten \cite{gaiottowitten} 研究过），与算术中的类似预测相符，尤其是
   Sakellaridis--Venkatesh \cite{sakellaridisvenkatesh} 的相对朗兰兹纲领。这一联系在
   Ben-Zvi--Sakellaridis--Venkatesh 正在进行的工作中得到发展。
\end{re}

	\ssszh 我们来到最后一点，先略微非正式地陈述。在代数簇上使用构造层时，有几类不同的层可以
	选择，它们表现相似。显然，人们通常使用其中之一：（i）D-模，其相应的上同调理论（即「常层」的整体
	截面）是代数 de Rham 上同调；（ii）$\ell$-进层，其相应的上同调理论是 étale 上同调；（iii）在复数上
	工作时，使用解析拓扑下的构造层，其相应的上同调理论是奇异上同调，即 Betti 上同调。

	在我们讨论的（量子）几何朗兰兹版本中，对应的构造部分属于 de Rham 型。也就是说，在自守侧我们研究
	$G$-丛模空间上的 D-模，在谱侧研究带平坦联络的 $\cG$-丛模空间（即 de Rham 意义下的局部系统）上的
	凝聚层。与此相关，对 $k$ 上光滑射影曲线 $X$，Arinkin--Gaitsgory \cite{arinkingaitssingsupp} 用现代
	方式表述的整体猜想断言等价
	%
	\[
	\Dmod(\Bun_G(X)) \simeq \ICn(\LS(X)).
	\]
	%
	正如我们现在要描述的，本理论也有 étale 与 Betti 风味的版本。

	\ssszh[受限变化的局部几何朗兰兹] 局部算术朗兰兹的近期重要进展，是把它在表示的族中表述，
	即讨论整个范畴，而不只是一个一个不可约表示
	\cite{emertonhelm}、\cite{hellmann}、\cite{zhufamilies}、\cite{farguescholze}。也就是说，其现代表述
	大致是：$G(\Fqt)$ 的全部表示与某个适当的 Galois 表示模空间上的凝聚层之间的等价。

	\begin{re} 事实上，本着 Zelevinsky、Lusztig 与 Vogan \cite{voganloclang} 工作的精神，自守侧更
	微妙，它还要与其它群（例如 $G$ 的某些内形式）的表示范畴粘合起来。
	\end{re}

    在系列工作 \cite{restrictedvar1}、\cite{restrictedvar2}、\cite{restrictedvar3} 中，
    Arinkin--Gaitsgory--Kazhdan--Raskin--Rozenblyum--Varshavsky
    给出了一个在特征零与函数域情形下都说得通的整体范畴几何朗兰兹猜想。其形式是
    %
    \begin{equation} \label{e:etalelang}
           \on{Shv}^{\text{étale}}_{\on{nilp}}(\on{Bun}_G(X)) \simeq
           \on{IndCoh}_{\on{nilp}}(\on{LocSys}_{\cG}^{\on{restr}}(X)).
    \end{equation}
%
	在左端，考虑的是整体曲线 $X$ 上 $G$-丛模空间上的 ind-构造 $\ell$-进层，其奇异支撑落在
	$T^* \Bun_G(X)$ 的整体幂幺锥中。在右端，遇到的是 $X$ 上 $\ell$-进 $\cG$-局部系统的某个适当定义的
	模空间上、带幂幺奇异支撑的归纳凝聚层。

	\begin{re}
	非常粗略地说，基本图景是：在复数上，谱侧看起来像是 $\cG$-局部系统的 Betti 模叠的「撕开」版本，
	即参数化拓扑基本群表示的模空间。也就是说，按局部系统的半单化把 Betti 叠分解为许多块，
	再取它们形式邻域的不交并。例如，不可约局部系统的形式邻域在 $\on{LocSys}_{\cG}^{\on{restr}}$ 中是孤立的，
	并等同于 Betti 或 de Rham 模空间中相应的形式完备化。与此相关，上标「restr」代表受限变化。
	\end{re}

	在函数域设定下，取 Frobenius 的范畴迹，可以证明这恢复出（紧支撑）自守形式，特别地，
	恢复出函数域上尖点表示的非分歧整体朗兰兹的自守到 Galois 方向。

	在未发表的工作中，Gaitsgory 表述了 \eqref{e:etalelang} 的局部版本。它大致断言等价
	%
	\[
	     \on{Shv}^{\text{étale}}(G(\overline{\mathbb{F}}_q(\!(t)\!))\mod \simeq
	     \TwoIC(\LS^{\on{restr}}),
	\]
	%
	其中左端实际使用的是更微妙的范畴表示概念，以修正全忠实嵌入
	\[
\on{Shv}^{\text{étale}}(X) \otimes \on{Shv}^{\text{étale}}(Y) \hookrightarrow \on{Shv}^{\text{étale}}(X \times Y)
\]
	并非本质满的问题，并进一步过渡到「谱有限」对象的满子范畴。

	通过 Frobenius 的范畴迹，预期在函数域设定下这能恢复 Zhu \cite{zhufamilies} 所表述的算术局部
	朗兰兹猜想的某种形式。

	\ssszh[局部 Betti 几何朗兰兹] 最后，Ben-Zvi--Nadler \cite{benzvinadlerbetti} 表述了关于 Betti 层的
	几何朗兰兹版本。
	此时，对复数上的射影代数曲线 $X$，整体等价取形式
	%
	\[
	  \on{Shv}^{\on{Betti}}_{\on{nilp}}(\Bun_G(X)) \simeq
	  \ICn(\LS^{\on{Betti}}(X) ).
	\]
	%
	在左端，考虑的是 $\Bun_G(X)$ 的解析化上的 Betti 层（不必是 ind-构造的），其奇异支撑落在整体
	幂幺锥中。在右端，考虑的是局部系统的 Betti 模空间上的归纳凝聚层，即把 $X$ 的底拓扑空间的特征簇
	自然地增强为导出叠。

	局部情形没有很好的文献记载，但我们可以概述一些基本部分。首先，正则奇点的部分由等价
	%
	\begin{equation} \label{e:bezfam}
	      \on{Shv}^{\on{Betti}}_{\on{nilp}}( \mathring{I}\backslash LG /
	      \mathring{I}) \simeq \ICn( \widetilde{\cG} / \cG \underset {\cG /
	      \cG} \times \widetilde{\cG} / \cG).
	\end{equation}
	%
	确立。非正式地说，这是 Bezrukavnikov 等价 \eqref{e:2realz} 及其「非幂幺单值、且单值在族中变化」
	的变体的一个版本，其中单值在伴随商 $\check{G} / \check{G}$ 中变化。\footnote{这一等价可能尚未见于
	文献。不过综述作者与 H. Chen 有一个论证，希望不久能记录成文。}

	按构造，\eqref{e:bezfam} 右端的模等同于 $\LS^{\on{Betti}}$，即 $\cG / \cG$ 上带幂幺奇异支撑的
	归纳凝聚范畴层。完整的断言（即包括野分歧的对应物）应当把环路群的 Betti 范畴表示（一阶近似下
	即 $\on{Shv}^{\on{Betti}}(LG)\mod$）与 Stokes 数据模空间上带幂幺奇异支撑的归纳凝聚范畴层联系起来。

	\ssszh 至此，本综述已到结尾。希望读者对本领域的一些基本想法、结果与猜想有所感受，也能欣赏它们的
	美。有兴趣的读者不妨深入我们巡览过的文献；在一个有如此多迷人开放问题的领域中，边做边学尤其可行。
